\chapter{Proposed Method -- Phương pháp đề xuất}\label{ch:method}
Chương này đặc tả phiên bản tham chiếu GeoI-BQ v1, rồi phân tích privacy và utility của cùng phiên
bản đó. Luồng chính là GPS $\rightarrow$ thông tin vị trí đã bảo vệ $\rightarrow$ belief $\rightarrow$
đoạn Q $\rightarrow$ dịch vụ POI $\rightarrow$ xử lý đáp án tại local. Các bước lý thuyết được phát
biểu trước benchmark để bảo đảm kết luận không phụ thuộc một attacker học cụ thể.

\section{Problem formulation and assumptions}\label{sec:formulation}
\subsection{Mạng đường, trajectory và nhu cầu dịch vụ}
Cho một mạng đường có hướng với tập trạng thái hữu hạn $V$. Mỗi trạng thái $v$ có tọa độ $a_v$ trong
một hệ chiếu mét cố định. Hai lane khác nhau có thể trùng tọa độ; chúng vẫn là hai trạng thái cho
việc chuyển tiếp. GPS thật $x_i$ không nhất thiết trùng đúng một $a_v$. Trong định lý Geo-I, khoảng
cách $d(x,x')=\|x-x'\|_2$ là khoảng cách Euclid của cùng phép chiếu. Chi phí di chuyển trên đường có
hướng được dùng riêng cho reachability và utility, không thay vào chứng minh tam giác mà không kiểm
tra các giả thiết tương ứng.

Trajectory $X=(x_0,x_1,\ldots)$ được định nghĩa trên mọi mốc \emph{có thể đọc} theo clock chung, kể cả
mốc bị bộ lọc bỏ qua do pacing/cap. Như vậy, hai đầu vào so sánh không bị đổi hệ chỉ số vì số lần đọc
thực tế khác nhau. Luận văn dùng metric trajectory
\begin{equation}
D_\infty(X,X')=\sup_i\|x_i-x_i'\|_2.
\label{eq:trajectory-metric}
\end{equation}
Với trace hữu hạn, supremum là maximum. Metric này đo độ lệch giữa hai giả thuyết trajectory tại
cùng clock; không được hiểu là khoảng cách mà một người đã di chuyển trong một chuyến.

Nhu cầu thật $\psi_t$ gồm purpose, loại POI và các tham số local như radius hoặc destination. Catalogue
POI là $\mathcal P$. Client cần trả một danh sách phù hợp đã sắp xếp theo GPS hiện tại và $\psi_t$.
Độ sâu đánh giá $k=5$ chỉ tạo reference top-five, không buộc ứng dụng chỉ trả năm địa điểm.

\subsection{Dữ liệu riêng, ngữ cảnh công khai và transcript}
Thông tin riêng gồm GPS, nhu cầu thật và các target có quan hệ với trajectory được khai báo trong
từng tác vụ. Ngữ cảnh $C$ gồm bản đồ, cấu hình, public clock, catalogue và conditional law của provider.
Lịch sử nội bộ đã bảo vệ $E$ gồm các nhánh noisy test và điểm REM. Belief được tính từ thông tin
này cùng bối cảnh cho phép; không đọc trực tiếp GPS để thay thư viện Q.

Transcript $O_n$ của server gồm các Q theo thứ tự, timestamp, request/reply và retry sau $n$ lần chọn
đoạn. GPS local, $\psi_t$, belief, ledger và cờ fresh/reuse không thuộc transcript. Đáp án đã lọc tại
thiết bị cũng không được gửi ngược lại trong mô hình này. Nếu một ứng dụng gửi click, destination
hoặc diagnostics riêng lên server, kênh quan sát đã thay đổi và phải được phân tích bổ sung.

Để chứng minh nhiều bước, lịch sử $\mathcal H_j$ giữ cả các đoạn đã chọn, nhãn lane giả và replies.
Nó mạnh hơn wire chỉ gồm tọa độ: lane aliases có thể cùng tọa độ nhưng dẫn tới thư viện khác ở bước
sau. Các nhãn lane đó thuộc hành động giả do mô hình sinh, không phải lane thật của người dùng.
Phân tích trên lịch sử mở rộng rồi chiếu xuống wire tránh giả định sai rằng tọa độ xác định duy nhất
toàn bộ trạng thái chuyển tiếp.

\subsection{Attacker và ranh giới bảo vệ}
Attacker là LSP hoặc bên có cùng quyền quan sát. Nó biết thuật toán, tham số công khai, mạng đường
và có thể lưu toàn transcript. Định lý không giả định attacker chỉ dùng một snapshot, chỉ biết kNN
hoặc không biết lịch phát thật. Hai traces $X,X'$ là phép so sánh của cơ chế; attacker không cần biết
sẵn hai trajectory đó trong một cuộc tấn công cụ thể.

Tuyến thật và các target tương lai trong phân tích được cố định ngoại sinh trước phép chạy. Nếu Q
hoặc đáp án POI làm thay đổi hành động tiếp theo của người dùng, cần mô hình tương tác có feedback.
Một user identity hoặc future target chỉ nhận hệ quả từ định lý khi mối quan hệ target--input--output
thỏa mô hình đã phát biểu. Account/IP/biển số trực tiếp và raw side information từ người tương quan
không mặc định được bảo vệ.

\subsection{Giả thiết dùng trong các định lý}\label{sec:assumptions}
\begin{enumerate}[label=\textbf{A\arabic*.}]
\item \textbf{Support chung.} Tại cùng $C,\mathcal H_{j-1}$, thư viện $\mathcal A_j$, thứ tự action và
masses nền giống nhau, mọi mass dương. GPS, Z, belief và nhu cầu riêng không loại ứng viên. Có một
đoạn đứng yên khả thi khi cần fallback.
\item \textbf{Upstream Geo-I lý tưởng.} REM và noisy test dùng các kernel xác suất được đặc tả;
downstream chỉ đọc $E$ và bối cảnh cho phép. Không có đường trực tiếp GPS $\rightarrow$ support,
request hoặc quyết định gửi riêng. Nhánh không gian qua restart cần phục hồi đầy đủ history, chi phí
và RNG state; khớp timestamp/config chưa đủ xác nhận việc phục hồi đó.
\item \textbf{Kernel Q.} Điểm vào kernel nằm trong $[0,1]$. Sampler integer sử dụng bit đều mới,
độc lập với private input và upstream randomness. Sai số belief không phá cận Q nếu score vẫn bị chặn.
\item \textbf{Ledger và lịch phát.} Commit trước gửi; counter/allowance không rollback hoặc reset
trong cùng phiên. Các khung còn lại và retry dùng đoạn đã chọn. Clock/mở/đóng phiên được giữ cùng
nhau khi so hai input, hoặc dừng theo quy tắc chỉ dùng lịch sử công khai.
\item \textbf{Provider law.} Tại cùng Q/history, server có cùng conditional response law cho hai
input. Cache suy từ replies; thông tin riêng không được đưa thêm vào kênh request/reply.
\item \textbf{Utility.} Metadata dùng chứng minh cache còn hợp lệ; client dùng cùng GPS, nhu cầu và
total order khi so các tập ứng viên. Muốn dùng floor từ reply signatures, batch nhận được phải đầy đủ
và khớp signatures. Calibration/proxy error chưa biết được giữ là chưa chứng nhận.
\end{enumerate}

Một phiên có cùng token/ledger. Restart tiếp tục phiên đó; mở nhiều phiên liên kết cần composition
riêng. Những giả thiết này được dùng có chọn lọc: cận Q có thể đúng dù upstream không thỏa Geo-I,
nhưng cận không gian và minimum của hai cận cần nhánh upstream hợp lệ.

\section{System architecture -- Vị trí mô hình trong hệ thống}
Hình~\ref{fig:system} phân biệt thiết bị tin cậy và LBS server. Cảm biến cung cấp GPS cho lớp bảo vệ
và cho local ranking. Privacy/planner tạo các tọa độ Q để lấy dữ liệu; server trả POI theo request
chung. Thiết bị giữ nhu cầu thật để chọn đáp án từ những metadata đã nhận.

\begin{figure}[htbp]
\centering\resizebox{\linewidth}{!}{\begin{tikzpicture}[>=Latex,b/.style={draw=ink!65,rounded corners=2pt,fill=white,text width=3.2cm,align=center,inner sep=7pt,font=\small},a/.style={->,thick,draw=ink},private/.style={->,thick,draw=teal}]
\node[b] (user) at (0,0) {\textbf{Người dùng / cảm biến}\\GPS $x_t$, nhu cầu $\psi_t$};
\node[b,draw=teal,fill=teal!5] (model) at (5.4,0) {\textbf{GeoI-BQ}\\Bảo vệ / lập kế hoạch Q};
\node[b] (server) at (10.8,0) {\textbf{LBS server}\\Nhận Q; trả POI};
\node[b] (context) at (5.4,-2.3) {\textbf{Ngữ cảnh công khai}\\Bản đồ, cấu hình, POI};
\node[b] (local) at (5.4,-4.8) {\textbf{Xử lý đáp án local}\\Union/cache, lọc, sắp xếp};
\node[b] (out) at (0,-4.8) {\textbf{Đáp án người dùng}\\Danh sách phù hợp\\đã sắp xếp};
\node[b] (attacker) at (10.8,-3) {\textbf{Attacker tại server}\\Lưu Q, clock, replies\\Suy luận trajectory};
\draw[a] (user)--node[above,text width=1.6cm,align=center,font=\scriptsize] {GPS cho lớp\\bảo vệ}(model);
\draw[a] (context)--(model);\draw[a] (model)--node[above,text width=1.6cm,align=center,font=\scriptsize] {Q, request\\chung}(server);
\draw[a] (server.south west)--(8.3,-1.5)--(8.3,-4.8)--node[above,font=\scriptsize] {POI}(local.east);
\draw[a] (local)--(out);\draw[a,dashed] (server)--(attacker);
\draw[private] (user.west)--(-2.2,0)--(-2.2,-3.5)--(5.4,-3.5)--node[right,font=\scriptsize] {GPS + $\psi$ local}(local.north);
\begin{scope}[on background layer]\node[fit=(user)(model)(context)(local)(out),draw=teal,fill=teal!2,rounded corners,inner sep=12pt] {};\end{scope}
\node[font=\small\bfseries,text=teal,fill=white,inner sep=2pt] at (2.7,1.05) {THIẾT BỊ TIN CẬY};
\end{tikzpicture}
}
\caption{Kiến trúc hệ thống và ranh giới quan sát. Chỉ Q/request chung đi sang LBS; GPS và nhu cầu thật phục vụ client.}
\label{fig:system}
\end{figure}

Z không phải một Q mặc định được phát lên server. Nó là trạng thái đã bảo vệ giúp quản lý private
reads và tạo belief. Q là hành động dịch vụ công khai. Nếu chỉ công bố Z, vẫn có thể có một bảo đảm
Geo-I lý tưởng nhưng chất lượng lấy POI bị gắn với một query location. Nếu chọn cả tập Q trực tiếp
từ GPS, cũng có thể xây dựng cơ chế đúng; khi đó phải phân tích trực tiếp toàn joint kernel và các
lần lặp theo thời gian. Hai tầng ở đây giúp phân vai và kiểm soát các điều kiện thiết kế.

\section{Model architecture and workflow}\label{sec:model-workflow}
Hình~\ref{fig:model} tổ chức phương pháp thành năm lớp vận hành. Trong mỗi lớp có các thao tác logic
riêng: bảo vệ/cập nhật vị trí, ước lượng trạng thái, tạo thư viện và chọn đoạn, phát khung, rồi hợp và
xử lý POI. Server không quan sát belief hoặc quyết định fresh/reuse. Đường phản hồi từ cache tới
planner chỉ mang thông tin suy từ replies, không đưa $\psi_t$ vào support hoặc cadence.

\begin{figure}[htbp]
\centering\resizebox{\linewidth}{!}{\begin{tikzpicture}[>=Latex,
 b/.style={draw=ink!65,rounded corners=2pt,fill=white,text width=8cm,align=center,inner sep=6pt,font=\small},
 e/.style={draw=blue!60,fill=blue!4,rounded corners,text width=3.1cm,align=center,inner sep=5pt,font=\small},
 a/.style={->,thick,draw=ink}]
\node[b] (rem) {\textbf{1. REM + noisy reuse + cap Geo-I}\\Tạo Z trên đường hoặc giữ Z cũ;\\dự toán chi phí trước GPS};
\node[e,text width=7.8cm,above=.85cm of rem] (gps) {\textbf{Input: GPS $x_t$ tại mốc được phép đọc}};
\node[b,below=.5cm of rem] (belief) {\textbf{2. Belief từ protected history}\\Dự đoán di chuyển; cập nhật\\bằng likelihood của Z};
\node[b,draw=teal,below=.5cm of belief] (select) {\textbf{3. Chọn chung một đoạn Q}\\Chấm điểm bằng belief; centered kernel bị chặn\\Thư viện và masses nền chỉ dùng public history};
\node[b,below=.5cm of select] (send) {\textbf{4. Phát ba khung đã commit}\\0/20/40 s, mỗi khung K5; L60 mỗi loại/Q};
\node[b,below=.5cm of send] (local) {\textbf{5. Union/cache và local ranking}\\GPS hiện tại + nhu cầu thật $\psi_t$};
\node[e,text width=7.8cm,below=.7cm of local] (answer) {\textbf{Output: toàn danh sách phù hợp đã sắp xếp}};
\node[e,right=1.1cm of belief] (pub) {\textbf{Public context}\\Đường, clock, Q trước, replies/cache};
\node[e,right=1.1cm of select] (ledger) {\textbf{Ledger Q}\\$\Gamma_Q$, counter, đoạn đã commit};
\node[e,right=1.1cm of send] (lbs) {\textbf{LBS server}\\Q $\rightarrow$ POI};
\draw[a] (gps)--(rem);\draw[a] (rem)--(belief);\draw[a] (belief)--(select);
\draw[a] (select)--(send);\draw[a] (send)--(local);\draw[a] (local)--(answer);
\draw[a,dashed] (pub.west)--++(-.5,0)|-(select.east);
\draw[a,dashed] (ledger.west)--(select.east);
\draw[a,blue] (send.east)--(lbs.west);\draw[a,blue] (lbs.south)|-(local.east);
\draw[a,teal] (gps.west)--++(-1.1,0)|-(local.west);
\node[rotate=90,font=\scriptsize,text=teal] at ($(belief.west)+(-1.3,-1.0)$) {GPS hiện tại cho local ranking};
\draw[a,dashed] (local.west)--++(-.45,0)|-(select.west);
\begin{scope}[on background layer]
\node[fit=(rem)(belief)(select)(send)(local),draw=teal,fill=teal!2,rounded corners,inner sep=10pt] (model) {};
\end{scope}
\node[anchor=south west,font=\scriptsize\bfseries,text=teal,fill=white,inner sep=2pt] at (model.north west) {MÔ HÌNH GEOI-BQ TẠI CLIENT};
\end{tikzpicture}
}
\caption{Kiến trúc nội bộ GeoI-BQ. Một đoạn ba khung được chọn chung; từng khung và retry sau đó là hậu xử lý.}
\label{fig:model}
\end{figure}

\subsection{Vấn đề, quyết định thiết kế và điều cần chứng minh}
REM sử dụng miền đường công khai để tránh đầu ra vô nghĩa ngoài mạng. Điều cần chứng minh là phân
phối trên miền đó thỏa Geo-I, không chỉ rằng một điểm được chọn có vẻ hợp lý. Noisy reuse tránh phải
tạo Z mới ở mọi lần phát; phép kiểm tra vẫn dùng GPS nên phải được tính vào cap.

Belief mô tả những vị trí có thể của người dùng để phục vụ retrieval. Nó không là chứng chỉ attacker
không suy ra được người dùng. Thư viện đoạn Q giữ độc lập với private pruning để có common support;
belief chỉ điều chỉnh xác suất trong một giới hạn. Đây là phần tách yêu cầu utility khỏi điều kiện
support của chứng minh.

Kernel centered và ledger Q giới hạn mức điều chỉnh theo thông tin riêng qua nhiều đoạn. Commit cả
đoạn bảo đảm các khung 20/40 s không lén trở thành những lần chọn mới chưa hạch toán. Common category
cover giữ payload purpose/category thật ở local. Union/cache và ranking dùng GPS để cải thiện đáp
án mà không tạo một lần phát GPS mới. Mỗi thiết kế sẽ được nối với mệnh đề tương ứng ở
Mục~\ref{sec:privacy} và~\ref{sec:utility}.

\subsection{Clock, bộ nhớ và cấu hình tham chiếu}
Một đoạn kéo dài 60 giây với các khung ở offset 0, 20, 40 s. Lần chọn đoạn tiếp theo ở offset 60 s.
Lớp bảo vệ dùng GPS mới tối thiểu cách nhau 60 s và chỉ khi dự toán cap cho phép. GPS local phục vụ
xếp hạng có thể được đọc theo nhu cầu ứng dụng; không suy từ pacing bảo vệ rằng tổng số GNSS reads
hoặc năng lượng đã giảm.

\begin{figure}[htbp]\centering
\begin{tikzpicture}[>=Latex,x=1cm,y=1cm,font=\small]
\draw[->] (0,0)--(13,0) node[right] {$t$};
\foreach \x/\t in {0/0,2/20,4/40,6/60,8/80,10/100,12/120}{
\draw (\x,.08)--(\x,-.08);\node[below] at (\x,-.1){\t};
\fill[blue] (\x,.65) circle (2pt);}
\node[anchor=east] at (-.3,.65){Gửi Q};
\foreach \x in {0,6,12}{\fill[teal] (\x,1.45) circle (2.4pt);}
\node[anchor=east] at (-.3,1.45){Có thể đọc};
\draw[teal,thick,decorate,decoration={brace,amplitude=5pt}] (0,2.1)--(5.85,2.1);
\draw[teal,thick,decorate,decoration={brace,amplitude=5pt}] (6.05,2.1)--(11.85,2.1);
\node at (3,2.55){Đoạn 1: $\varepsilon_{Q,1}=1/2$};
\node at (9,2.55){Đoạn 2: $\varepsilon_{Q,2}=1/6$};
\end{tikzpicture}
\caption{Hai clock khác nhau. Dấu ``có thể đọc'' còn phụ thuộc dự toán cap; phát khung không đồng nghĩa lấy mẫu Q hoặc đọc bảo vệ GPS mới.}
\label{fig:clock}
\end{figure}

Cấu hình tham chiếu dùng $K=5$, $L=60$ mỗi loại/tọa độ, sáu categories công khai và reference $k=5$.
Vì vậy một packet có tối đa $5\times6\times60$ reply entries trước loại trùng, không phải chừng đó
POI unique. Prototype có catalogue 418 POI, nên union có thể bão hòa; đây là điều kiện cần khi giải
thích Recall cao. Planner sử dụng portfolio beam 24, novelty/long-range weights 4 và travel scale
60 s. Các tham số là cấu hình đã đóng băng, không là tối ưu phổ quát được chứng minh.

\section{Algorithm specification -- Đặc tả thuật toán}\label{sec:algorithm}
\subsection{REM trên miền đường công khai}
Với $u>0$ và cùng $V$, cơ chế lấy mẫu trạng thái và ánh xạ sang tọa độ:
\begin{align}
R_u(v\mid x)&=\frac{\exp[-u\|x-a_v\|_2/2]}{\sum_{w\in V}\exp[-u\|x-a_w\|_2/2]},\label{eq:rem}\\
V_t&\sim R_u(\cdot\mid x_t),\qquad Z_t=a_{V_t}.\nonumber
\end{align}
Các trạng thái gần GPS có trọng số lớn hơn, nhưng mọi trạng thái trong miền chung có xác suất dương
ở kernel lý tưởng. Không lọc từng điểm theo một điều kiện ``đã đạt Geo-I'' rồi random đều. Geo-I là
tính chất của \emph{toàn phân phối} khi thay input. Cũng không có bước planar Laplace rồi mới map vào
REM trong đặc tả này. Hệ số $1/2$ chia độ lệch giữa tử số và normalizer trong chứng minh.

\subsection{Noisy reuse và cap từng phiên}
Sau Z đầu, tại mốc được phép đọc, với tham chiếu cũ $h$, đặt
\begin{equation}
D_t=\mathbf1\{\|x_t-h\|_2+\xi_t>\theta\},\qquad
\xi_t\sim\operatorname{Lap}(0,1/u),\quad \theta=200\text{ m}.
\label{eq:private-test}
\end{equation}
$D_t=0$ giữ $Z_t=h$; $D_t=1$ gọi REM để tạo Z mới. Tạo đầu tiên tốn một đơn vị; test rồi giữ tốn
một; test rồi tạo mới tốn hai. Bộ lọc dự toán chi phí lớn nhất cần cho thao tác \emph{trước khi đọc}
GPS. Nhánh không đọc không được dùng một phép kiểm tra GPS riêng để quyết định rằng nó sẽ không đọc.

Các số ngân sách có quan hệ:
\begin{equation}
u=\frac{B}{2H_{\rm cfg}},\qquad M=2H_{\rm cfg}-1,\qquad C_s=Mu.
\label{eq:budget-config}
\end{equation}
Trong prototype, $B=0{,}24$ m$^{-1}$, $H_{\rm cfg}=12$, nên $u=0{,}01$ m$^{-1}$, $M=23$ và
$C_s=0{,}23$ m$^{-1}$. $H_{\rm cfg}$ là tham số accounting, không đòi biết trước route hoặc thời
lượng chuyến. Cap không đặt thời gian tối đa người dùng được đi; hết allowance thì lớp bảo vệ tiếp
tục từ thông tin đã bảo vệ, đổi lại khả năng thích nghi với GPS mới giảm.

Muốn một chính sách rủi ro không gian có tỷ số tối đa $R_*$ trong bán kính so sánh $d_*$, điều kiện đủ
của cận thô là $C_s\le\log(R_*)/d_*$. Công thức cho phép suy ra $u=C_s/M$ rồi kiểm tra utility.
Luận văn không khẳng định cấu hình 0,23 đã được chọn từ một yêu cầu rủi ro như vậy ngay từ đầu, hoặc
là một bán kính nhiễu. Đó là cap hiệu lực của cấu hình tham chiếu.

\subsection{Belief từ lịch sử đã bảo vệ}
Trên các trạng thái đại diện $s$, bắt đầu bằng prior công khai. Dự đoán và cập nhật:
\begin{align}
b_t^-(s)&=\sum_r b_{t-1}(r)T_{\Delta t}(r,s),\label{eq:belief-prediction}\\
b_t(s)&=\frac{b_t^-(s)\ell_t(Z_t\mid s,h)}{\sum_r b_t^-(r)\ell_t(Z_t\mid r,h)}.\label{eq:belief}
\end{align}
Ở lần đầu, $\ell$ cộng REM probabilities của các lane có cùng tọa độ. Sau đó, likelihood của quan
sát Z bao gồm cả xác suất reuse và refresh. Nếu $r_s$ là xác suất giữ ở trạng thái đại diện $s$:
\begin{equation}
\ell(z\mid s,h)=r_s\mathbf1_{\{z=h\}}+(1-r_s)\sum_{v:a_v=z}R_u(v\mid s).
\label{eq:emission}
\end{equation}
Vì REM có thể chọn lại cùng tọa độ cũ, không suy fresh/reuse chỉ từ việc hai Z có bằng nhau không.
Mốc không có private read chỉ dự đoán; không coi Z cũ lặp lại là một quan sát REM độc lập mới.

Lưới, prior và chuyển tiếp là xấp xỉ công khai. Prototype dùng một đại diện gần tâm cho mỗi cell
belief và một mô hình chuyển tiếp theo thời gian; nó không được gọi là posterior đã hiệu chuẩn của
người thật. Privacy của kernel Q chỉ cần score bị chặn. Sai số belief có ý nghĩa riêng ở phần chuyển
sang utility thực, nơi không được tự đặt sai số bằng 0.

\subsection{Thư viện đoạn, điểm số và lấy mẫu integer}
Một action $A_a$ gồm ba khung $Q_a^{(0)},Q_a^{(1)},Q_a^{(2)}$, mỗi khung có $K$ slots. Đường chuyển
tiếp của các slot phải khả thi trên mạng có hướng từ fake state trước. Public search/portfolio và
fallback đứng yên tạo thư viện $\mathcal A_j(\mathcal H_{j-1})$. Không dùng GPS/Z/belief để cắt bỏ
action. Điều kiện floor public nếu không đạt phải được báo giảm chứng chỉ hoặc dùng fallback chung,
không chuyển sang một support riêng theo người dùng.

Với bảng proxy công khai bị chặn $g_a(s)$ và belief $b$, tính raw score $v_a$. Sau chuẩn hóa và
lượng tử hóa thành integer $F_a$, đặt $z_a=F_a/S$, $S=65536$:
\begin{align}
v_a&=\sum_s b(s)g_a(s),\qquad 0\le z_a\le1,\qquad \mu=\sum_a w_a z_a,\nonumber\\
\alpha_j&=\frac{\varepsilon_{Q,j}}{2+\varepsilon_{Q,j}},\qquad
p_a=w_a\bigl[1+\alpha_j(z_a-\mu)\bigr].\label{eq:kernel}
\end{align}
$w$ là phân phối nền chung với masses dương. Belief điều chỉnh probabilities, không chọn support.
Min--max trước quantization có thể làm thay đổi thang đo; $z_a$ không là Recall và không mặc định
lệch raw $v_a$ tối đa $1/S$. Những cận không dựa vào sensitivity của phép min--max này được chứng
minh trực tiếp từ $z\in[0,1]$.

Allowance của lần chọn thứ $j$ là
\begin{equation}
\varepsilon_{Q,j}=\frac{\Gamma_Q}{j(j+1)},\qquad\Gamma_Q=1.
\label{eq:q-allowance}
\end{equation}
Lịch công khai này không đòi biết trước số đoạn. Viết $w_a=m_a/T$, $\alpha_j=N/D$, với các integer
dương thích hợp. Kernel được lấy mẫu qua trọng số integer:
\begin{equation}
W_a=m_a\left[DTS+N\left(TF_a-\sum_c m_cF_c\right)\right],\qquad T=\sum_c m_c.
\label{eq:integer-kernel}
\end{equation}
Chọn một integer đều trên $[0,\sum_aW_a)$ rồi dò cumulative weights cho đúng PMF hữu hạn.
Điều này tách sai số tính score trước lượng tử hóa khỏi việc lấy mẫu PMF, nhưng vẫn cần nguồn bit
đều mới; một PRNG dùng trong mô phỏng không tự là chứng chỉ entropy thông tin lý thuyết.

\subsection{Phát khung, hợp POI và local ranking}
Sau khi lấy mẫu, commit cả đoạn, counter và ledger trước khi gửi khung đầu. Các khung 20/40 s dùng
đoạn đã lưu; retry dùng lại bản tin. Request hỏi mọi loại POI với cùng L60 ở mỗi Q. Server không
nhận purpose/radius/destination thật qua giao diện này.

Client hợp các metadata nhận được, loại trùng theo ID POI và giữ cache theo chính sách đã khai báo.
Tại thời điểm dùng dịch vụ, GPS hiện tại và nhu cầu thật tạo bộ lọc/total order local. Danh sách
được sắp xếp chỉ gồm các POI thực nhận và còn hợp lệ; local ranking không khôi phục metadata chưa
được truy hồi. Catalogue tĩnh trong benchmark giúp cache tích lũy; trường hợp hết hạn/status động
đòi kiểm tra lại giả thiết utility.

\subsection{Giả mã và chi phí vận hành}
\begin{enumerate}
\item Đọc public clock, kiểm tra pacing và dự toán Geo-I cap. Quyết định có thể đọc GPS phải dựa
trên trạng thái/lịch cho phép trước private read.
\item Nếu được đọc: tạo REM đầu tiên hoặc noisy test rồi reuse/refresh; cập nhật $E$ và chi phí.
Nếu không: tiếp tục từ thông tin đã bảo vệ.
\item Dự đoán belief theo thời gian; cập nhật emission chỉ khi có private read.
\item Ở đầu đoạn: tạo thư viện/masses nền chung, tính $v,z,p$ và lấy mẫu bằng integer kernel;
commit action và allowance lần $j$.
\item Phát khung đã commit; nhận replies, union/cache hợp lệ. Không dành thêm allowance cho các
khung còn lại hoặc retry cùng bản tin.
\item Lọc/sắp xếp tại local bằng $x_t,\psi_t$ và trả toàn danh sách phù hợp.
\end{enumerate}

Một REM trực tiếp duyệt $|V|$ weights; belief và planner có chi phí theo số trạng thái đại diện,
transition và candidate segments. Public portfolio có thể dùng cache/tables nhưng không vì vậy
mà suy ra latency của điện thoại. Khi score POI đã có, hợp các reply entries là tuyến tính theo số
entries và sắp xếp $|U_t|$ ứng viên mất $O(|U_t|\log|U_t|)$ phép so sánh. Chi phí truy vấn đường và
cập nhật metadata phụ thuộc provider/context; số liệu thời gian ở Chương~\ref{ch:experiments} ghi rõ
phạm vi đo, đặc biệt output replay không đo lại generation của proposed.

\subsection{Ví dụ từ belief tới một đoạn Q}
Xét ba trạng thái mock ở 0, 100, 300 m với prior $(0{,}5;0{,}3;0{,}2)$ và $u=0{,}01$. Nếu Z đầu
ở 100 m, likelihood REM của quan sát đó theo ba GPS giả thuyết xấp xỉ $(0{,}3315;0{,}5065;0{,}2312)$.
Nhân prior và chuẩn hóa cho belief $(0{,}4554;0{,}4175;0{,}1271)$.

Hai đoạn có proxy trên ba trạng thái là $(0{,}9;0{,}6;0{,}2)$ và $(0{,}3;0{,}7;0{,}9)$, cho raw scores
khoảng $0{,}686$ và $0{,}543$. Sau min--max, $z=(1;0)$. Với base $w=(1/2;1/2)$ và allowance đầu
$\varepsilon_{Q,1}=1/2$, ta có $\alpha=1/5$, $\mu=1/2$, nên $p=(0{,}55;0{,}45)$.
Đoạn hữu ích hơn được ưu tiên nhưng đoạn kia vẫn có probability dương. Đây là ví dụ tính toán,
không thay cho chứng minh hoặc điểm benchmark.

\subsection{Ví dụ nhiều bước từ output đã lưu}
Chọn cố định family TEST đầu theo ID, \path{cache-029}, trip/draw đầu, không chọn theo score.
Bảng~\ref{tab:sample} hiển thị GPS tương đối so với lúc bắt đầu, Z nội bộ, Q1 và số POI unique trong
cache. Q2--Q5 vẫn có trong packet; tọa độ trong bảng làm tròn, artifact giữ đủ precision.
\Needspace{210pt}
\par\begingroup\footnotesize\setstretch{1.04}\setlength{\parskip}{0pt}
\begin{longtable}{@{}P{.6cm}P{2.3cm}P{2.7cm}P{3.3cm}P{3.3cm}P{2.394cm}@{}}
\caption{Một prefix nhiều bước cố định từ family cache-029}\label{tab:sample}\\
\toprule
\textbf{s} & \textbf{$\Delta$GPS (m)} & \textbf{Nhánh / units tích lũy} & \textbf{Z (lat; lon)} & \textbf{Q1 (lat; lon)} & \textbf{POI cache} \\
\midrule\endfirsthead
\multicolumn{6}{l}{\small\textit{Bảng \thetable\ (tiếp)}}\\\toprule
\textbf{s} & \textbf{$\Delta$GPS (m)} & \textbf{Nhánh / units tích lũy} & \textbf{Z (lat; lon)} & \textbf{Q1 (lat; lon)} & \textbf{POI cache} \\
\midrule\endhead
\midrule\multicolumn{6}{r}{\small\textit{Còn tiếp ở trang sau}}\\\endfoot
\bottomrule\endlastfoot
0 & (0; 0) & Tạo Z / 1 & (39{,}9662; 116{,}3280) & (39{,}9771; 116{,}3588) & 381 \\
20 & (120; 3) & Không đọc / 1 & (39{,}9662; 116{,}3280) & (39{,}9771; 116{,}3588) & 381 \\
40 & (236; 26) & Không đọc / 1 & (39{,}9662; 116{,}3280) & (39{,}9771; 116{,}3588) & 381 \\
60 & (355; 30) & Tạo Z / 3 & (39{,}9664; 116{,}3331) & (39{,}9770; 116{,}3586) & 385 \\
120 & (711; 18) & Tạo Z / 5 & (39{,}9662; 116{,}3322) & (39{,}9764; 116{,}3534) & 402 \\
180 & (412; 33) & Giữ Z / 6 & (39{,}9662; 116{,}3322) & (39{,}9759; 116{,}3529) & 407 \\
\end{longtable}\endgroup\par


Ở 20/40 s, mô hình phát các khung còn lại mà không đọc bảo vệ GPS mới. 60/120 s tạo Z mới với tổng
chi phí tương ứng; 180 s giữ Z nhưng có thể chọn một đoạn mới từ belief đã cập nhật. Units nhân
$u=0{,}01$ cho chi phí Geo-I tích lũy. Z, GPS và nhánh trong bảng là diagnostics để giải thích, không
được đưa vào attacker features. Cache tăng trong mẫu không phải định lý rằng Recall luôn tăng theo
thời gian khi người dùng di chuyển.
